\documentclass[10pt,a4paper]{amsart}
\usepackage[margin=19mm]{geometry}
\usepackage{amsmath,amssymb,mathtools}
\usepackage[scr=rsfs,cal=euler]{mathalfa}
\usepackage{xfrac}
\usepackage[unicode,hidelinks,bookmarks=false]{hyperref}
\newcommand{\ie}{i.e.\ }
\newcommand{\cf}{cf.\ }
\newcommand{\resp}{respectively\ }
\newcommand{\Subsection}{\subsection}
\providecommand{\nobreakdash}{\nobreak\hbox{-}}
\setlength{\emergencystretch}{2em}
\multlinegap0pt
\usepackage{mathtools}
\newtheorem{proposition}{Proposition}
\usepackage{cleveref}
\crefname{proposition}{Proposition}{Propositions}
\title{Binary X-rays of doubly stochastic matrices}
\author{Bangzheng Li, Yuewei Liu, Yuan Yao}
\date{Source: arXiv:2608.01442v2 (CC BY 4.0)}
\begin{document}
\maketitle

\noindent\emph{Source and licence.} Binary X-rays of doubly stochastic matrices. Version arXiv:2608.01442v2 (CC BY 4.0), distributed by arXiv under the Creative Commons Attribution 4.0 International licence, http://creativecommons.org/licenses/by/4.0/, as recorded in the arXiv licence metadata for this exact version and shown on the abstract page https://arxiv.org/abs/2608.01442v2. Full licence notice: \texttt{licenses/arxiv-2608.01442-CC-BY-4.0.txt}.\par\smallskip

\noindent\textit{Editorial scope.} Counted unit: the article's proposition prop:x-conditions (source0.tex lines 78--82). The article's complete original proof of it is delivered with it (source0.tex lines 101--119, one \texttt{proof} environment of 19 lines, which the article places away from the statement). No mathematics is written, repaired or re-derived: the statement and the proof are the article's own lines, in the article's own notation, and no retained line is substituted or re-flowed.  No further source-local context is needed: every cross-reference of every retained line resolves inside the two delivered spans.  Nothing was omitted from any retained span, so the package carries no omission record.  No retained line contains a citation, so the wrapper carries no bibliography and no bibliography entry.  No retained line contains a figure environment, an image-inclusion command or drawing code, so no image is delivered and the package declares no asset dependency.  Editorial formatting only, in addition: an AMS \texttt{amsart} preamble that restates the article's own declarations and macros used by the retained lines, namely the environments \texttt{proposition} on separate counters, together with the standard packages those lines need.  The retained environments print in this document as Proposition 1 in order of appearance, whereas the article prints the same items with the numbering of its own full document.  The complete native source of the article as distributed in the arXiv e-print for this exact version is bundled unchanged as \texttt{sources/206628/source0.tex}.
\begin{proposition}\label{prop:x-conditions}
    An integer vector $s\in \{-n+1, \dots, n-1\}^n$ satisfies the permutohedral constraints (i.e.\ is a score vector) if and only if its corresponding frequency vector $x\in \mathbb{N}^{2n-1}$ satisfies the equation $S: \sum_{d=-n+1}^{n-1} x_d = n$ and the following inequalities for all integers $k = 0, \dots, n$: 
        \[D_k: \sum_{i=1}^{2k} (2k-i+1) x_{-n+i} \leq k(k+1) \quad \text{and} \quad U_k: \sum_{i=1}^{2k} (2k-i+1) x_{n-i}\leq k(k+1).\]
    (Here we let $x_{-n} = x_n = 0$.)
\end{proposition}

\begin{proof}[Proof of \cref{prop:x-conditions}]
    We first consider a score vector $s$ and a tournament with that score vector. If $s$ is a permutation of $(-n+1, -n+3, \dots, n-3, n-1)$, then the corresponding tournament is ``totally ordered'', where a higher-ranked player always wins over a lower-ranked player. In this case the corresponding frequency vector is $(1, 0, 1, \dots, 1, 0, 1)$, and it is not difficult to see that all inequalities in the statement are in fact equalities in this case. 
    
    Otherwise, it is not difficult to see that there exist two players $i$ and $j$ where $s_i \geq s_j$ but $i$ did not win over $j$. If we change the outcome of this match to be 1 point more in favor of $i$ (a draw into a win, or a loss into a draw), then the sum of squares of all scores strictly increases (since $(s_i+1)^2 + (s_j-1)^2 \geq s_i^2 + s_j^2 + 2$), so this process must terminate in finitely many steps, ending in a totally ordered tournament. By running this process in reverse, we see that any frequency vector of a score vector can be obtained via starting with $(1, 0, 1, \dots, 1, 0, 1)$ and repeatedly decrementing $x_i$ and $x_j$ by $1$ for some pair $(i,j)$ where $j - i \geq 2$ and $x_i, x_j\geq 1$ while incrementing $x_{i+1}$ and $x_{j-1}$ by $1$. It is clear that such operation never makes any of the listed inequalities false, so the final vector $x$ must also satisfy them.

    For the reverse direction, we first claim that the permutohedral constraints are equivalent to the following (quadratic) inequalities for all $t = 1, \dots, 2n-1$: \begin{align*}
        \sum_{i=1}^{t}(-n+i)x_{-n+i} &\geq -\left(\sum_{i=1}^{t}x_{-n+i}\right)\left(n-\sum_{i=1}^{t}x_{-n+i}\right), \\
        \sum_{i=1}^{t}(n-i)x_{n-i} &\leq \left(\sum_{i=1}^{t}x_{n-i}\right)\left(n-\sum_{i=1}^{t}x_{n-i}\right),
    \end{align*}
    or equivalently 
    \begin{equation}\label{eq2}
        \left(\sum_{i=1}^{t}x_{-n+i}\right)^2 \leq \sum_{i=1}^t ix_{-n+i}\quad \text{and}\quad \left(\sum_{i=1}^{t}x_{n-i}\right)^2 \leq \sum_{i=1}^t ix_{n-i}.
    \end{equation}
    These inequalities follow from the permutohedral constraints by letting $I$ be the set of all indices $i$ for which $s_i\leq -n+t$ (or $s_i\geq n-t$). 
    
    We can also show that these inequalities are sufficient. First, assume without loss of generality that the S-vector is sorted, so it suffices to prove the permutohedral constraints for $I = [k]$ (and $I = [n]\setminus [k]$) for each $k \in [n-1]$. Suppose that $I = [k]$ and $s_k = -n+t$, and that there are $x'\leq x_{-n+t}$ copies of $-n+t$ in the first $k$ coordinates. Since the first inequality in \eqref{eq2} is convex with respect to the variable $x_{-n+t}$, and it holds with $x_{-n+t}$ replaced by $0$ (which is the same as the inequality for $t-1$), it must also hold for with $x_{-n+t}$ replaced by $x'$ (since $0\leq x' \leq x_{-n+t}$). This is equivalent to the permutohedral constraint for $I = [k]$. The case for $I = [n]\setminus [k]$ is symmetric.

    Now, for each $t$ let $k = \sum_{i=1}^t x_{-n+i}$ (which is an integer in $[0, n]$). Since $(2k-i+1)x_{-n+i}$ is non-negative when $i\leq k$ and non-positive when $i > k$, by $D_k$ we have \[(2k+1)k-\sum_{i=1}^t ix_{-n+i} = \sum_{i=1}^{t} (2k-i+1) x_{-n+i} \leq \sum_{i=1}^{k} (2k-i+1) x_{-n+i}\leq k(k+1),\] which directly rearranges to the first inequality of \eqref{eq2}. The second inequality can be proved symmetrically.
\end{proof}

\end{document}
